%% 00-abstract.tex

The Europeana Data Model (EDM) enables structural interoperability across
heterogeneous cultural heritage institutions, but its design trades semantic
precision for breadth: the same property (\texttt{dc:type}) carries
controlled-vocabulary terms in library records, free-text strings in museum
records, and structural hierarchy identifiers in archive records.
Schema-level ontology alignment tools, which operate on lexical and structural
similarity between ontology definitions, cannot recover these sector-specific
conventions from the specification alone.
Corpus-based analysis is therefore a precondition for reliable automated
alignment of heterogeneous EDM metadata.

We present a corpus-driven three-layer alignment methodology for mapping
DDB-EDM records to domain ontologies (RDA, RiC-O, VRA Core, Music Ontology),
mediated through \textsc{mocho}, a mid-level hub ontology for cross-domain
cultural heritage objects.
The three layers address distinct alignment challenges: (1) property alignment
via a data-driven DC\,$\to$\,RDA sub-property mapping table derived from
corpus profiling; (2) archival hierarchy typing via a lookup from DDB
\texttt{hierarchyType} codes to RiC-O and DoCO classes; and (3) document-type
dispatch via a three-level fallback table keyed on \texttt{dc:type} value,
institutional sector, and mediatype, assigning WEMI-level classes from domain
ontologies.
Applied to a cross-domain pilot corpus of 115{,}432 DDB records spanning all
seven institutional sectors, the pipeline produces 44.4 million RDF triples.
All scripts and dispatch tables are publicly available \todo{confirm URL}.
